\section{The timestep confound}
\label{sec:confound}

Section~\ref{sec:worth} left AWWL's MRG positive but not significant after
Holm correction. This section examines one thing AWWL alters besides the
spectrum, the timestep weighting that every noise-dependent band weighting
carries (Proposition~\ref{prop:confound}). Trained without the other, AWWL's
timestep schedule and its frequency shape at constant total weight are each
significantly worse than MSE, and the two partly cancel when combined. The
spectral change comes mainly from the shape, whereas how AWWL's gap to MSE
divides between the two depends on the order in which they are added.

\subsection{AWWL's weighting as a timestep schedule}
\label{sec:confound-schedule}

For an orthogonal transform, the per-sample loss of Eq.~\eqref{eq:family} is
the pixel mean square $\operatorname{ms}(r)$ weighted by the total
$S(\sigma) = \sum_b \omega_b(\sigma)$, a timestep weight that does not depend
on the residual, plus a zero-sum band term that carries the frequency shape
(Proposition~\ref{prop:confound}(a)). For AWWL at the CIFAR-10 setting
($\alpha = 0.2$, $p = 1$), $S$ (Eq.~\eqref{eq:total}) falls from nearly
$1 - \alpha = 0.8$ at the lowest noise level to $\alpha = 0.2$ at the highest
(Fig.~\ref{fig:weights}(a)). Its mean over the $T = \num{1000}$ training
timesteps, \wTotalMean, lies near the high-noise value because
\pctSigmaHigh\% of them have $\sigma_t > 0.9$; under $\sigma$ uniform on
$[0,1]$ the mean is exactly $1/2$ for every $\alpha$ and $p$, because
$\lambda_p(1-\sigma) = 1 - \lambda_p(\sigma)$. This residual-independent
schedule runs opposite to Min-SNR-$\gamma$, which damps the low-noise
timesteps: over the training timesteps, the Pearson correlation of
$S(\sigma_t)$ with the Min-SNR multiplier $\kappa_t$
(Section~\ref{sec:notation}) is $\rSMinSNRCifar$ at this setting and
$+\rSMinSNRDB$ at the DreamBooth setting ($\alpha = 0.8$, $p = 2$;
Fig.~\ref{fig:weights}(b)). By Proposition~\ref{prop:confound}(c) the sign is
that of $2\alpha - 1$, so the two published optima, presented as balances
between the approximation and detail bands, also impose timestep schedules of
opposite direction.

\guardedfigure{0.85}{weights.png}{%
  AWWL's weighting as a timestep schedule. (a) The weights of
  Eq.~\eqref{eq:weights} at the CIFAR-10 setting ($\alpha = 0.2$, $p = 1$):
  $w_{\mathrm{LL}}$, the detail weight $w_{\det}$ (applied to the mean of the
  three detail bands) and their total $S$ (Eq.~\eqref{eq:total}); the dotted
  line marks the mean of $S$ over the training timesteps; the top axis gives
  $\sigma_t$. (b) The total at the CIFAR-10 and DreamBooth settings
  ($\alpha = 0.8$, $p = 2$, Stable Diffusion's scaled-linear schedule) and the
  Min-SNR-$\gamma$ multiplier $\kappa_t$ ($\gamma = 5$) on each schedule, each
  divided by its mean over the training timesteps; the legend gives the
  Pearson correlation $r$ of each total with $\kappa_t$.}{fig:weights}

\subsection{Separating the schedule from the shape}
\label{sec:confound-controls}
\label{sec:confound-results}

Two normalisation controls keep AWWL's band shape and hold one
residual-independent measure of its per-timestep weight constant. The first
rescales the weights to a constant total,
$\tilde\omega_b = c\,\omega_b / S(\sigma)$, which removes $S$ exactly
(Proposition~\ref{prop:confound}(d)); it uses $c = \levelSumFixed$, the
uniform-$\sigma$ mean of $S$, and is repeated at $c = 1$, which changes only
the effect of AdamW's stability constant $\varepsilon_{\mathrm{A}}$ and of
gradient clipping (Proposition~\ref{prop:scale}(a)). The second fixes $\sum_b \omega_b^2$ at
\levelSqFixed, AWWL's mean over the training timesteps. For band gradients of
equal norm, the two normalised quantities bound the per-timestep gradient
weight: fixing $S$ fixes it when the gradients are parallel, fixing
$\sum_b \omega_b^2$ when they are orthogonal, as in output space, and neither
control is the stronger (Proposition~\ref{prop:scale}(b);
Appendix~\ref{app:gradnorm}). A third arm, the timestep-only arm, keeps the
schedule and removes the shape: it trains $S(\sigma_t)\operatorname{ms}(r)$,
the first term of Eq.~\eqref{eq:shape-decomp}, and its two contrasts, with
MSE and with AWWL, form a Holm family of their own.

% Generated by scripts/make_tables.py --- do not edit by hand.
% Regenerate after every sweep; the paper reads these directly.

\begin{table}[tbp]
\centering
\footnotesize
\setlength{\tabcolsep}{3pt}
\caption{The timestep schedule and the normalisation controls of AWWL (released form, detail bands averaged), CIFAR-10, epoch 199, DDIM-100. The band weights are $\omega_b = (w_{\mathrm{LL}}, w_{\det}/3, w_{\det}/3, w_{\det}/3)$, and $S(\sigma_t) = \sum_b \omega_b = w_{\mathrm{LL}} + w_{\det}$ is AWWL's total weight. The timestep-only arm trains the pixel MSE multiplied by $S(\sigma_t)$, AWWL's timestep component alone (Proposition~\ref{prop:confound}); each normalisation control multiplies the AWWL loss by a function of $\sigma_t$ alone and holds the stated sum fixed over $\sigma_t$, the sum of squares at AWWL's mean over the training timesteps, $\sum_b \omega_b^2 = 0.053$. $\Delta$ is the paired-by-seed FID difference, the arm minus MSE or minus AWWL, with its 95\% $t$-interval over the $n$ seeds; $p$ is the two-sided paired $t$-test, raw and Holm-corrected. Holm families: against MSE, the 21-arm family of Table~\ref{tab:compare-fid}; against AWWL, the four confound contrasts of this table and Table~\ref{tab:confound-summed}; $^{\dagger}$for the timestep-only arm, its own two-test family (this arm against MSE, and AWWL against this arm), corrected separately from both. Bold: $p_{\mathrm{Holm}} < 0.05$ within the stated family. Lower panel: the difference from MSE of four of the arms, arranged by whether $S(\sigma_t)$ is AWWL's or held constant and whether AWWL's frequency shape is present. Only $S$ is held constant: the effective timestep weight $\Lambda(t)$ is not (Propositions~\ref{prop:confound}(b) and~\ref{prop:residual}). The sum-of-squares control holds $\sum_b \omega_b^2$ rather than $S$ and has no cell.}
\label{tab:confound}
\begin{tabular}{>{\raggedright\arraybackslash\hangindent=1em\hangafter=1}p{3.3cm}rrrrrrr}
\toprule
 & & \multicolumn{3}{c}{Against MSE} & \multicolumn{3}{c}{Against AWWL} \\
\cmidrule(lr){3-5}\cmidrule(lr){6-8}
Arm & $n$ & $\Delta$ [95\% CI] & $p$ & $p_{\mathrm{Holm}}$ & $\Delta$ [95\% CI] & $p$ & $p_{\mathrm{Holm}}$ \\
\midrule
\multicolumn{8}{l}{\emph{Without a frequency shape}} \\
MSE & 5 & \emph{reference} &  &  &  &  &  \\
AWWL, timestep only & 5 & $+0.81$\;$[0.48, 1.14]$ & $0.0025$ & $\mathbf{0.0050}$$^{\dagger}$ & $+0.37$\;$[0.004, 0.73]$ & $0.048$ & $\mathbf{0.048}$$^{\dagger}$ \\
\addlinespace[2pt]
\multicolumn{8}{l}{\emph{With AWWL's frequency shape}} \\
AWWL & 5 & $+0.44$\;$[-0.15, 1.03]$ & $0.11$ & $0.32$ & \emph{reference} &  &  \\
AWWL, $\sum_b \omega_b = 1$ & 5 & $+1.15$\;$[0.48, 1.82]$ & $0.0090$ & $\mathbf{0.045}$ & $+0.71$\;$[0.35, 1.07]$ & $0.0056$ & $\mathbf{0.012}$ \\
AWWL, $\sum_b \omega_b = 0.5$ & 5 & $+1.34$\;$[0.98, 1.70]$ & $0.0005$ & $\mathbf{0.0049}$ & $+0.90$\;$[0.48, 1.31]$ & $0.0039$ & $\mathbf{0.012}$ \\
AWWL, $\sum_b \omega_b^2$ fixed & 5 & $+0.88$\;$[0.47, 1.30]$ & $0.0041$ & $\mathbf{0.024}$ & $+0.44$\;$[0.23, 0.65]$ & $0.0041$ & $\mathbf{0.012}$ \\
\bottomrule
\end{tabular}
\par\medskip
\begin{tabular}{@{}l>{\raggedright\arraybackslash}p{4.1cm}>{\raggedright\arraybackslash}p{4.1cm}@{}}
\toprule
 & \multicolumn{2}{c}{Frequency shape} \\
\cmidrule(lr){2-3}
Timestep total $S(\sigma_t)$ & absent & AWWL's \\
\midrule
held constant & MSE\newline $0$ (reference) & AWWL, $\sum_b \omega_b = 0.5$\newline $+1.34$\;$[0.98, 1.70]$ \\
\addlinespace[2pt]
AWWL's & AWWL, timestep only$^{\dagger}$\newline $+0.81$\;$[0.48, 1.14]$ & AWWL\newline $+0.44$\;$[-0.15, 1.03]$ \\
\bottomrule
\end{tabular}
\end{table}


Both normalisation controls are worse than MSE
($p_{\mathrm{Holm}} = \dSumFixedHolm$ and \dSqFixedHolm{} over the
\nFamilyMain{} arms) and worse than AWWL, by \dSumFixedVsAWWL{} and
\dSqFixedVsAWWL~FID ($p_{\mathrm{Holm}} = \dSumFixedVsAWWLHolm$ and
\dSqFixedVsAWWLHolm{} over the \nFamilyConfound{} confound contrasts;
Table~\ref{tab:confound}); the run at $c = 1$ repeats the first on the same
seeds and is not an independent replication. On the summed-detail shape of
Eq.~\eqref{eq:summed}, fixing $\sum_b \omega_b^2$ improves on the
unnormalised parent, so the sign of the effect depends on the band shape
(Appendix~\ref{app:supp-confound}).

With MSE and AWWL, these arms form a two-by-two design over whether $S$ varies
with $\sigma_t$ and whether the frequency shape is present
(Table~\ref{tab:confound}, lower panel). Against MSE
($S$ constant, no shape), the schedule alone costs $+\dTimestep$~FID
$[\dTimestepLo, \dTimestepHi]$ ($p_{\mathrm{Holm}} = \dTimestepHolm$ in its
two-test family) and the shape at constant $S$ costs $+\dSumFixed$
($+\dSqFixed$ with $\sum_b \omega_b^2$ held instead), yet AWWL, which
combines them, is not resolved from MSE ($+\dAWWL$, raw $p = \dAWWLP$). The
components therefore partly cancel, and the design reads from either corner:
adding the shape to the schedule changes FID by $\dAWWLVsTimestep$
$[\dAWWLVsTimestepLo, \dAWWLVsTimestepHi]$ (raw and Holm
$p = \dAWWLVsTimestepP$, at the edge of significance), and adding the
schedule to the shape at constant $S$ changes it by $-\dSumFixedVsAWWL$
($p_{\mathrm{Holm}} = \dSumFixedVsAWWLHolm$). The schedule produces
little of the spectral change: its coefficient-weighted high-band change is
\corrHighWTimestep~dB $[\corrHighWTimestepLo, \corrHighWTimestepHi]$, against
\corrHighWAWWL~dB for AWWL and \corrHighWSumFixed{} and \corrHighWSqFixed~dB
for the two controls, and is worth \valueTimestep~FID
$[\valueTimestepLo, \valueTimestepHi]$, against \priceTransplantValue{} for
AWWL. AWWL's spectral change therefore comes mainly from its shape, but its
cost does not divide so simply: the schedule alone has an MRG of
\mrgTimestep~FID $[\mrgTimestepLo, \mrgTimestepHi]$ (raw $p = \mrgTimestepP$;
the arm is outside the priced-arm family), close to AWWL's \mrgMain, whereas
the shape at constant $S$ has larger MRGs (\mrgArmSumFixed{} with
$\sum_b \omega_b = \levelSumFixed$, \mrgArmSqFixed{} with $\sum_b \omega_b^2$
fixed; Table~\ref{tab:pricing-arms}), so the design credits AWWL's residual
cost to neither component alone.

This interaction is described, not tested, and the rows of the design are
defined by $S$, not by the effective timestep weight $\Lambda(t)$, the ratio of
the expected loss to the expected pixel mean square at $t$
(Eq.~\eqref{eq:eff-weight}); it equals $1$ for MSE and
$S(\sigma_t)$ for the timestep-only arm, but for the controls, as for AWWL,
it retains a residual-dependent term that no rescaling by a function of
$\sigma$ removes (Proposition~\ref{prop:confound}(b),~(d)) and that changes
with $\sigma_t$ under a Gaussian data model (Proposition~\ref{prop:residual}).
The `constant' row is therefore constant in $S$ only, and as $\Lambda(t)$ was
not measured, the difference between the columns cannot, in either row, be
assigned to the frequency shape alone.

\subsection{The confound across the objective family}
\label{sec:confound-family}

The residual-dependent part of the timestep weight is present throughout the
family. By Proposition~\ref{prop:confound}(b), $\Lambda(t)$ is a convex
combination of the pixel-equivalent weights $(N/n_b)\,\omega_b(\sigma_t)$
with the residual's band energy fractions as coefficients; unless those
weights are equal, it depends on the residual's spectrum, which changes with
$\sigma_t$, whatever normalisation fixes $S$. FFL, whose weights contain no
$\sigma_t$, weights timesteps through the ratio $\Lambda_{\mathrm{FFL}}(r)$ of
its loss to the pixel MSE, which depends on the spectral shape of the residual
(Section~\ref{sec:family-ffl}). This reweighting was not measured, so FFL's
result (Section~\ref{sec:worth}) is not divided into timestep and frequency
parts.
Whether an objective reweights timesteps must therefore be established from
$\Lambda(t)$, not from whether $\sigma_t$ appears in its definition.
